% Canonical node 013 · M005; current section path 1.2.6.
% Canonical owner: M005
\cnnaNodeHeading{013}{M005}{Conductance-unit normalization independence}
\cnnaNodePlacement{1.2.6}{D2}

\paragraph*{Formal statement.}
N001 chooses $C_\star=1$ as the canonical bootstrap unit.  M005 defines equality
of normalized scalar responses without division by
\begin{equation}
  (R,C)\sim(R',C')
  \quad\Longleftrightarrow\quad
  RC'=R'C.
  \label{eq:m005-cross-product}
\end{equation}
For every rational scale $\lambda>0$,
\begin{equation}
  (R,C)\sim(\lambda R,\lambda C),
  \label{eq:m005-rescaling}
\end{equation}
and for every positive rational unit $u$ and normalized value $x$,
\begin{equation}
  (x,1)\sim(ux,u).
  \label{eq:m005-canonical-representative}
\end{equation}

\paragraph*{Proof strategy and interpretation.}
Equation~\eqref{eq:m005-rescaling} is the associativity and commutativity identity
$R(\lambda C)=(\lambda R)C$.  Positivity is not needed for that algebraic
identity; it identifies the rescaling as an admissible change of a positive
conductance unit.
% CNNA-EXTREF-BEGIN EXT-USE-M005-SCALING-DIRECT
\cite[Electrical network preliminaries]{DoyleSnell1984ElectricNetworks}.
% CNNA-EXTREF-END EXT-USE-M005-SCALING-DIRECT
Python evaluates the corresponding quotient exactly with
\texttt{Fraction} and rejects nonpositive units and scales.  Lean uses
\texttt{Rat} and the cross-product relation, avoiding division and all
regularization.

\paragraph*{Necessity checks.}
Scaling only $R$ or only $C$ generally changes the normalized value.  A zero
scale satisfies the bare cross-product identity but cannot represent a unit,
and a negative scale violates the positive-conductance convention even though
it preserves an algebraic ratio.  These distinctions prevent the theorem from
being misread as an unrestricted gauge symmetry.

% CNNA-INFOBOX-BEGIN INFO-M005-SCALE-ORBIT
\begin{cnnaInfoBox}{Positive scaling classes, not a gauge field}
The positive pairs $(R,C)$ and $(\lambda R,\lambda C)$ lie on one orbit of the
multiplicative action of $\mathbb Q_{>0}$, and $R/C$ is constant on that orbit.
Choosing $C_\star=1$ selects a convenient representative.
% CNNA-EXTREF-BEGIN EXT-USE-M005-SCALE-INFOBOX
\cite[Sec.~I, current-balance and reduced conductance equations]{DoerflerBullo2013KronReduction}.
% CNNA-EXTREF-END EXT-USE-M005-SCALE-INFOBOX
No local gauge
group, gauge connection, edge-dependent rescaling symmetry or dynamical gauge
field is introduced at this node.
\end{cnnaInfoBox}
% CNNA-INFOBOX-END INFO-M005-SCALE-ORBIT

\paragraph*{Result, limits and handoff.}
M005 proves that the fixed N001 numeral is a unit representative rather than a
third physical input.  The theorem is scalar and rational and makes no claim of
homogeneity for the full Schur/DtN map.  Verified edge \texttt{E013} certifies
C014.  Verified edge \texttt{E068} supplies the fixed representative
$C_\star=1$ consumed by M003 after the C015 identity transform.

\noindent\textbf{Primary code anchors.}
Python exact quotient and rescaling checks: lines 19--46; test: lines 14--25.
Lean cross-product relation, N001 rational lift and two rescaling theorems:
lines 19--73.  The Supplement lists all nine anchors and exact hashes.
